\chapter{Introduction}
\label{ch:introduction}
% MOTIVATION
\section{Motivation}
Since the launch of Sputnik 1 in 1957, the amount of human-made orbital debris has increased dramatically.
The first two pieces of debris were the rocket stage that launched the artificial satellite and the satellite itself.
Today, there are over 27,000 officially cataloged objects in orbit, most of which are non-functional and non-cooperative.  % https://www.nasa.gov/mission_pages/station/news/orbital_debris.html
Such non-cooperative systems in orbit have the potential of causing a self-sustaining growth of collisions with other space debris---known as the Kessler syndrome---that could make some low Earth orbits unusable for long-term use.
% These collisions have involved operational satellites in the past such as the collision of Iridium 33, a communications satellite, with Cosmos 2251, an inactive Russian military communications satellite in 2009.
Our research is inspired by the need to control and deflect tumbling non-cooperative objects in orbit so as to avoid the risk of collisions with more vital orbiting systems such as functional satellites and space stations.

% LITERATURE REVIEW
\section{Literature Review}
Non-cooperative space debris is generally inert and shut down in orbit.
Such spacecraft may actively or passively not provide information regarding their inertial and orbital parameters as a result of which external observation is vital.
Techniques for rendezvous with and control of such tumbling objects have been presented using a single spacecraft agent.
This thesis relates to the work done in \cite{albee2021robust} on autonomous rendezvous and docking with a tumbling target in orbit.
This paper proposes a pipeline involving the estimation of the tumble of the target, the prediction of the target state, the trajectory planning of the chaser, the correction of reference trajectory error (due to inertial estimate errors) using robust control, and docking with the target.
Similarly, \cite{li2019extended} designs a robust output-feedback control strategy for spacecraft rendezvous and docking maneuvers in the presence of external disturbances and actuator saturation.
The problem of fuel-optimal trajectory planning for rendezvous with tumbling targets has been addressed from the practical perspective of trajectory discretization and curve fitting \cite{sternberg2018parameterization}, hence exploring the need to balance computational efficiency with fuel cost.
The target satellite is assumed to be uncommunicative, rigid, without flexible structures, or fuel slosh, and the chaser satellite is taken to have accurate state estimation and control, both of which are assumptions considered in this thesis as well.
However, the paper assumes the synchronicity of the chaser with the target due to the short duration of the approach, hence naturally decreasing the likelihood of a collision between the chaser and the target.
We relax this assumption in this thesis as a result of which the chaser uses model predictive control with robust techniques to recursively plan paths that prevent collisions with the target.

Although this has been conceptualized using single-agent systems, we also explore the idea of using multiple chasers in order to bring a large tumbling target to a standstill.
For instance, target information estimation has been proposed \cite{nakka2022information} using multiple chasers and the subsequent state estimation, guidance, and control of the chasers.
Similarly, previous literature \cite{mercier2023relative} proposes multi-agent inspection of the target and the subsequent centralized estimation of agent positions with no a priori information available using sliding window factor graphs (SWFG).
However, these papers focus on the observation of the target and do not propose techniques for interacting with the target for subsequent detumbling and redirection into a graveyard orbit.
We explore the idea of extending this concept into multi-agent docking and interaction with the target.
Previous studies have explored multi-agent detumbling methods utilizing non-contact approaches such as electromagnetic satellites equipped with high-temperature superconducting (HTS) coils \cite{yu2021prospects}.
However, this thesis delves into the feasibility of employing the propulsion system of the satellite, which is already in place for orbital positioning, to achieve detumbling of the target object.

On a related note, previous literature \cite{nakka2022information} presents multi-agent solutions to on-orbit inspection from a distance.
This has been supplemented in \cite{choi2023resilient} by resilient techniques to optimally perform multi-agent inspection in the scenario of failures/attacks of one or more inspector chasers.
On a similar note, \cite{aurand2023exposure} uses decentralized deep reinforcement learning techniques to decide viewpoints, plan agent trajectories, and perform viewpoint transfers for the inspection of a target object using one or multiple inspector satellites.
We intend to explore the idea of blending the idea of using multiple chasers with the problem of detumbling the target.

%% Accuracy of parameter estimation in orbital missions
It is vital to estimate the position, velocity, and other inertial parameters of the target in the approach of the chaser to the target.
Techniques to estimate the inclination of the orbit and the rate of change in the right ascension of the ascending node and, hence, the right ascension of the ascending node (RAAN) have been proposed \cite{hanada2022estimation} with the assumption of a circular orbit.
Similarly, \cite{ge2020motion} achieves convergence for parameter estimation of a non-cooperative target in orbit using a stereo-vision-based extended Kalman filter (EKF) observer, however, assuming the orbit to be circular with no drag.
Now, previous research that deals with the problem of this thesis such as \cite{albee2021robust} and \cite{starek2017fast} assumes the accurate calculation of potentially unknown inertial and orbital parameters.
In an effort to relax this assumption, we explore the idea of robustness in trajectory planning using tube model predictive control.
In particular, we incorporate the concept of dynamic robust control-invariant (RCI) tubes as proposed in \cite{brett2019dynamic}.
This is further observed in application in \cite{specht2023autonomous}.

%% How can chasers interact physically with a target?
Chaser-target interaction can be made possible using several techniques.
These include the use of tethers by the chaser to attach onto a target that is at a distance \cite{zhang2019trajectory,olivieri2020test,liu2021trajectory}.
While an operational satellite or spacecraft can be either partially cooperative (via communication of target characteristics and dynamics) in the case of the on-orbit servicing problem, much of space debris is non-cooperative and passive \cite{albee2021robust}.
In addition, the tumble of the target may not allow the attachment of a free-flying tether at a specific location on the target.
Previous research \cite{park2017nonlinear} has dealt with this by instead proposing techniques to perform rendezvous and docking with the target.
Although this is a feasible solution, the possibility of collisions is very likely.
The problem of collisions has been addressed in \cite{stoneman2016nonlinear} by considering convex polytope models for the chaser and the target, hence, making the collision constraints more conservative.
A similar problem is solved using conservative constraints around obstacles and control of error dynamics in \cite{ali2019safe} to generate safe trajectories for the chaser.
Furthermore, tube model predictive control has been incorporated in \cite{specht2023autonomous} to constrain the trajectory of the chaser in the presence of uncertainty, so as to prevent collisions.

%% On-Orbit Servicing
Previous research on on-orbit servicing has dealt with a problem similar to that proposed and solved in this thesis.
The need to inspect, relocate, restore, and augment existing systems in orbit using servicing vehicles has been proposed in the past \cite{long2007orbit}.
The system proposed by \cite{battistini2022iterative} utilizes measurements of onboard sensors on both the target and the chaser and information exchange (for range and bearing measurements) via an inter-satellite link (ISL).
Also, \cite{ogilvie2008autonomous} considers the servicing of the NextSat satellite using the ASTRO servicer where the target satellite could determine and control its attitude but did not have thrusters to modify its orbit.
Hence, much research on on-orbit servicing assumes the sharing of information between the controlled spacecraft; i.e., the chaser, and the uncontrolled spacecraft; i.e., the target such that the chaser obtains accurate estimates of the angular and orbital properties of the target.

%% Observation of Non-Cooperative Spacecraft
Existing literature \cite{li2022target} also provides insight into the detection of uncooperative spacecraft in orbit using lightweight convolutional neural network models.
Similarly, the reconstruction of non-cooperative targets has been proposed with the use of variants of the neural radiance field (NeRF) algorithm \cite{caruso20233d} - the Instant NeRF algorithm reduces computation cost by several orders of magnitude whereas the Dynamic NeRF algorithm considers the motion of the target for reconstruction, hence making it possible to observe the target under changing attitudes and lighting conditions.
Furthermore, previous research \cite{meng2019identification} has proposed the identification of the inertial parameters of the target such as its mass, center of mass, principal moments of inertia, and offset quaternion with respect to any given body-fixed reference frame.
The paper suggests the use of a time-of-flight camera to observe the target followed by the motion state estimation using an extended Kalman filter.
The mass and the real moments of inertia are calculated from observation of how the target reacts to a linear impulse provided by a physical touch from the observing spacecraft.

%% THESIS CONTRIBUTIONS
\section{Thesis Contributions}
Based on the literature review that has been conducted in light of the motivation for this problem and the subsequent need for reliable solutions to the detumbling problem in orbit, we propose the application of adaptive dynamic tube model predictive control (ADTMPC) to facilitate chaser trajectory planning for rendezvous and docking.
We propose trajectory planning of the chaser using model predictive control where changes in the environment are considered in the problem by optimizing a finite time horizon but only implementing the solution for a much smaller duration and then optimizing again, repeatedly.
However, this may still lead to the possibility of collisions with the target.
This is all the more likely since the target is the main obstacle in this problem while the final docking position also lies on the surface of the target.
Moreover, we consider the possibility of uncertainty in trajectory due to unknown inertial and orbital parameters.
As a result, we propose the use of robust techniques, particularly, robust tubes, to set a bound on the possible trajectory of the chaser with the assumption of a range of values for the unknown parameters.
While robust tubes prevent collisions, they also prevent the chaser from reaching the final state that lies on the target (which is also the obstacle under consideration).
Hence, we propose techniques based on \cite{morozov2020performance} to shrink these robust tubes with better estimates of the unknown parameters that cause this trajectory uncertainty.
This in turn makes the system less conservative for trajectory planning while remaining safe from collisions with the target.

%% THESIS OUTLINE
\section{Thesis Outline}
Chapter \ref{sec:relevant-knowledge} presents the problem statement in more detail while also providing information on the setup of the chaser and the target.
Chapter \ref{sec:chaser_ampc} deals with the basic trajectory optimization problem followed by the estimation of unknown parameters using adaptive control techniques.
It further elaborates on the constraints and bounds involved in this problem. 
This is followed by an explanation of the effects of using dynamic tubes to avoid constraint violation in Chapter \ref{sec:dynamic-tube}.
It also proposes techniques to expand this problem using multiple chasers.
The chaser, following its attachment onto the target, needs to exert torque on the target to stop detumbling.
This is discussed in Chapter \ref{sec:target-deflection}.
In Chapter \ref{sec:results}, we demonstrate the effectiveness of using ADTMPC over other techniques that do not include elements of this algorithm.
Finally, the appendices provide more information on the dynamics of the chaser for a simplistic (circular orbit with no drag) scenario followed by a more realistic (elliptical orbit with quadratic drag) one.